% Part: first-order-logic
% Chapter: completeness
% Section: completeness-theorem

\documentclass[../../../include/open-logic-section]{subfiles}

\begin{document}

\iftag{FOL}
      {\olfileid[hi]{fol}{com}{cth}}
      {\olfileid[hi]{pl}{com}{cth}}

\olsection{पूर्णता प्रमेय}

\begin{explain}
अब अपने परिणामों को एक साथ रखें: हम पूर्णता प्रमेय पर पहुँचते हैं।
\end{explain}

\begin{thm}[पूर्णता प्रमेय]
\ollabel{thm:completeness}
मान लें कि $\Gamma$, !!{sentence}s का समुच्चय है। यदि $\Gamma$ संगत है, तो
वह संतुष्टनीय है।
\end{thm}

\begin{proof}
मान लें कि $\Gamma$ संगत है।\iftag{FOL}{ \olref[hen]{lem:henkin} से कोई
  संतृप्त संगत समुच्चय $\Gamma' \supseteq \Gamma$ है।}{}
\olref[lin]{lem:lindenbaum} से कोई $\Gamma^* \supseteq \iftag{FOL}{\Gamma'}{\Gamma}$
है जो संगत और !!{complete} है।
\iftag{FOL}{
  चूँकि $\Gamma' \subseteq \Gamma^*$, प्रत्येक
  !!{formula}~$!A(x)$ के लिए $\Gamma^*$ में इस रूप का !!a{sentence} है:
  \iftag{prvEx}
        {$\lexists[x][!A(x)] \lif !A(c)$}
        {$\lnot\lforall[x][!A(x)] \lif \lnot !A(c)$}।
  इसलिए $\Gamma^*$ संतृप्त है। यदि $\Gamma$ में~$\eq$ नहीं है, तो}{तब}
\olref[mod]{lem:truth} से
$\iftag{FOL}{\Sat{M(\Gamma^*)}}{\pSat{v(\Gamma^*)}}{!A}$ तभी और केवल तभी जब
$!A \in
\Gamma^*$। इससे विशेषतः निकलता है कि सभी $!A \in
\Gamma$ के लिए $\iftag{FOL}{\Sat{M(\Gamma^*)}}{\pSat{v(\Gamma^*)}}{!A}$,
अतः $\Gamma$ संतुष्टनीय है।\iftag{FOL}{
  यदि $\Gamma$ में~$\eq$ है, तो \olref[ide]{lem:truth} से सभी
  !!{sentence}s~$!A$ के लिए
  $\Sat{\equivclass{M}{\approx}}{!A}$ तभी और केवल तभी जब $!A \in \Gamma^*$।
  विशेषतः, $\Sat{\equivclass{M}{\approx}}{!A}$ सभी $!A \in
  \Gamma$ के लिए, अतः $\Gamma$
  संतुष्टनीय है।}{}
\end{proof}

\begin{cor}[पूर्णता प्रमेय, दूसरा रूप]
\ollabel{cor:completeness}
सभी $\Gamma$ और !!{sentence}s~$!A$ के लिए: यदि $\Gamma \Entails !A$, तो
$\Gamma \Proves !A$।
\end{cor}

\begin{proof}
ध्यान दें कि \olref{cor:completeness} और \olref{thm:completeness} में आने
वाले $\Gamma$ सार्वत्रिक रूप से परिमाणित हैं। भ्रम से बचने के लिए
\olref{thm:completeness} को किसी दूसरे चर से फिर कहें: !!{sentence}s के किसी
भी समुच्चय~$\Delta$ के लिए, यदि $\Delta$ संगत है, तो वह संतुष्टनीय है।
प्रतिपक्ष से, यदि $\Delta$ संतुष्टनीय नहीं है, तो $\Delta$ असंगत है। हम इसी
से उपप्रमेय सिद्ध करेंगे।

मान लें कि $\Gamma \Entails !A$। तब \olref[syn][sem]{prop:entails-unsat} से
$\Gamma \cup \{\lnot !A\}$ असंतुष्टनीय है। $\Gamma
\cup \{\lnot !A\}$ को अपना $\Delta$ मानने पर
\olref{thm:completeness} का पिछला रूप देता है कि $\Gamma \cup \{\lnot !A\}$
असंगत है।
\iftag{FOL}{%
  \tagrefs{prfAX/{fol:axd:prv:prop:prov-incons},
    prfSC/{fol:seq:prv:prop:prov-incons},
    prfND/{fol:ntd:prv:prop:prov-incons},
    prfTab/{fol:tab:prv:prop:prov-incons}}}{%
  \tagrefs{prfAX/{pl:axd:prv:prop:prov-incons},
    prfSC/{pl:seq:prv:prop:prov-incons},
    prfND/{pl:ntd:prv:prop:prov-incons},
    prfTab/{pl:tab:prv:prop:prov-incons}}} से
$\Gamma \Proves !A$।
\end{proof}

\tagprob{FOL}
\begin{prob}
\olref[fol][com][cth]{cor:completeness} से
\olref[fol][com][cth]{thm:completeness} को सिद्ध कीजिए और इस प्रकार दिखाइए
कि पूर्णता प्रमेय के दोनों रूप समतुल्य हैं।
\end{prob}
\tagendprob

\tagprob{notFOL}
\begin{prob}
\olref[pl][com][cth]{cor:completeness} से
\olref[pl][com][cth]{thm:completeness} को सिद्ध कीजिए और इस प्रकार दिखाइए
कि पूर्णता प्रमेय के दोनों रूप समतुल्य हैं।
\end{prob}
\tagendprob

\tagprob{FOL}
\begin{prob}
किसी !!a{derivation} प्रणाली के पूर्ण होने के लिए उसके नियम इतने प्रबल होने
चाहिए कि वे प्रत्येक असंतुष्टनीय समुच्चय को असंगत सिद्ध कर सकें। पूर्णता को
सिद्ध करने के लिए !!{derivation} के कौन-से नियम आवश्यक थे? क्या इनमें से
कोई नियम प्रमाण में कहीं प्रयुक्त नहीं हुआ? इन प्रश्नों के उत्तर के लिए ऐसी
सूची या आरेख बनाइए जो दिखाए कि पूर्णता के प्रमाण तक ले जाने वाले किन
परिणामों में !!{derivation} के कौन-से नियम प्रयुक्त हुए। इन प्रमाणों में
नियमों के प्रत्येक निःशब्द उपयोग को भी अवश्य अंकित कीजिए।
\olref[fol][com][cth]{thm:completeness}।
\end{prob}
\tagendprob

\tagprob{notFOL}
\begin{prob}
किसी !!a{derivation} प्रणाली के पूर्ण होने के लिए उसके नियम इतने प्रबल होने
चाहिए कि वे प्रत्येक असंतुष्टनीय समुच्चय को असंगत सिद्ध कर सकें। पूर्णता को
सिद्ध करने के लिए !!{derivation} के कौन-से नियम आवश्यक थे? क्या इनमें से
कोई नियम प्रमाण में कहीं प्रयुक्त नहीं हुआ? इन प्रश्नों के उत्तर के लिए ऐसी
सूची या आरेख बनाइए जो दिखाए कि पूर्णता के प्रमाण तक ले जाने वाले किन
परिणामों में !!{derivation} के कौन-से नियम प्रयुक्त हुए। इन प्रमाणों में
नियमों के प्रत्येक निःशब्द उपयोग को भी अवश्य अंकित कीजिए।
\olref[pl][com][cth]{thm:completeness}।
\end{prob}
\tagendprob

\end{document}
